\documentclass[11pt]{article}
\usepackage[margin=0.9in]{geometry}
\usepackage{hyperref}
\usepackage{listings}
\usepackage{xcolor}
\usepackage{booktabs}
\usepackage{enumitem}
\usepackage{longtable}
\usepackage{array}
\usepackage{amsmath}
\usepackage{ragged2e}

% Lenient line breaking for long \texttt strings.
\emergencystretch=4em
\tolerance=9999
\hbadness=10000
\sloppy

\lstset{
  basicstyle=\ttfamily\small,
  breaklines=true,
  columns=fullflexible,
  frame=single,
  backgroundcolor=\color{gray!10},
  showstringspaces=false
}

\newcolumntype{L}[1]{>{\RaggedRight\arraybackslash\hangindent=1em}p{#1}}
\newcommand{\vc}{\textbf{very confident}}
\newcommand{\gs}{\textbf{guessed}}
\newcommand{\nc}{\textbf{will not guess}}
\newcommand{\es}{\textsc{essential}}
\newcommand{\nh}{\textsc{nice-to-have}}

\title{COARDS Audit of \texttt{gradients\_by\_period/} on the URI OPeNDAP Server\\
       \large grep-dap test case, Prompt \#3}
\author{Compiled for P. Cornillon}
\date{2026-05-16}

\begin{document}
\maketitle

\begin{abstract}
This document does for the \texttt{gradients\_by\_period/} branch
what Prompt~\#2 did for \texttt{SST\_Orbits/} --- with one major
asymmetry. The orbit branch has CF-1.5--labelled metadata and the
Prompt~\#2 audit was mostly a question of whether what is there is
COARDS-correct. The gradient branch, by contrast, has
\emph{zero} declared metadata: no global attributes, no variable
attributes, no coordinate variables, no fill values, no units, no
long names. Everything --- including the orientation of the lat
and lon axes --- has to be inferred from the variable-name
grammar, the data values themselves, and the companion branch.
The audit accordingly groups each proposed attribute on two axes:
\textbf{importance} (\es{} vs.\ \nh{}) and \textbf{confidence}
(\vc{} / \gs{} / \nc{}). The six resulting buckets organize the
rest of the document.
\end{abstract}

\tableofcontents

\section{Background}

The \texttt{gradients\_by\_period/} branch contains 270 NetCDF
files, one per calendar month from January 2003 through December
2024, named \texttt{augmented\_monthly\_stats\_MM\_YYYY.nc}. Each
file declares exactly two dimensions
(\texttt{lon=360}, \texttt{lat=180}) and 34 variables, all
\texttt{Float64} of shape (\texttt{lon}, \texttt{lat}); the schema
is identical across the period
(\texttt{scripts/gradients\_dmr\_inventory.py}).

The variable names follow a regular grammar:

\begin{lstlisting}
{day|night}_pixel_count                    # SST validity count
{day|night}_as_pixel_count                 # along-scan gradient count
{day|night}_at_pixel_count                 # along-track gradient count
{day|night}_sum_SST[_squared]              # sum (and sum^2) of SST
{day|night}_sum_eastward_gradient[_squared]
{day|night}_sum_northward_gradient[_squared]
{day|night}_sum_magnitude_gradient[_squared]
{day|night}_sum_grad_as_per_km[_squared]
{day|night}_sum_grad_at_per_km[_squared]
{day|night}_sum_grad_mag_per_km[_squared]
\end{lstlisting}

That naming, plus the companion branch \texttt{SST\_Orbits/} (which
\emph{is} CF-labelled and whose own
\texttt{eastward\_gradient}/\texttt{northward\_gradient} variables
have \texttt{units="C/km"} and
\texttt{standard\_name="*\_temperature\_gradient"}), is the source
of essentially every inference below.

\section{What the data themselves disclose}\label{sec:data}

The pre-audit probes establish the following facts that are needed
later. (Scripts:
\texttt{scripts/gradients\_geom\_resolve.py},
\texttt{scripts/gradients\_signs\_test.py}.)

\paragraph{Grid orientation.} The pixel-count field has a clear
land/ocean pattern. Indexed against six known-land points (Sahara,
Australia interior, Amazon, Central US, Greenland, Tibet) and six
known-ocean points (Pacific centre, Atlantic centre, Indian
Ocean, North Pacific, Persian Gulf, warm pool), only one of the
four candidate conventions classifies all 12 correctly:
\textbf{lat\_idx 0 = $-89.5^\circ$} (south--north), \textbf{lon\_idx
0 = $-179.5^\circ$} ($-180\to+180$, cell-centered). The
others place Australia under ocean indices, or Greenland under
ocean indices, etc.

\paragraph{Cell-centered $1^\circ$ grid.} Implied:
\texttt{lat[k] = -89.5 + k} for $k\in[0,179]$,
\texttt{lon[k] = -179.5 + k} for $k\in[0,359]$.

\paragraph{Signs.} \texttt{eastward\_gradient},
\texttt{northward\_gradient}, \texttt{grad\_as\_per\_km}, and
\texttt{grad\_at\_per\_km} are signed (positive and negative sums
observed). \texttt{magnitude\_gradient} and
\texttt{grad\_mag\_per\_km} are strictly non-negative.

\paragraph{Sum semantics.}
\texttt{sum\_magnitude\_gradient} per cell is on average
$\sim$12$\times$ larger than $\sqrt{S_e^2+S_n^2}$ from the same
cell, so the magnitude sum is the sum of per-pixel magnitudes,
\emph{not} the magnitude of the vector sum.

\paragraph{Pixel counts.} \texttt{day\_pixel\_count} (mean 10\,030
per cell) $\neq$ \texttt{day\_as\_pixel\_count} (mean 8\,110)
$\neq$ \texttt{day\_at\_pixel\_count} (mean 4\,090). The
$N_{\rm at} \approx \frac{1}{2}N_{\rm as}$ ratio is consistent with
MODIS Aqua's 10-line scan structure, in which every tenth
along-track derivative straddles a scan boundary and is masked.

\paragraph{Divisor for the gradient magnitude.} The per-cell mean
$\mathtt{sum\_grad\_mag\_per\_km}/\mathtt{at\_pixel\_count}$
$\approx 0.044\,^\circ$C/km closely tracks the
$\mathtt{sum\_magnitude\_gradient}/\mathtt{pixel\_count}$
$\approx 0.046\,^\circ$C/km, while using \texttt{as\_pixel\_count}
yields $\sim 0.022$. So
\texttt{sum\_grad\_mag\_per\_km} pairs with
\texttt{at\_pixel\_count} (the more restrictive of the two).

\paragraph{Time.} No time variable exists in the file. The single
calendar month each file represents must be derived from the
filename (\texttt{MM\_YYYY}); $\texttt{time}_\text{center}$ would
naturally be the mid-month UTC date as a CF-style
\texttt{days since 1970-01-01 00:00:00 UTC}.

\paragraph{Fill value.} Observed: NaN. Counts (\texttt{*\_pixel\_count})
contain $0$ explicitly (over land) and \texttt{NaN} where the
producer chose not to write a zero, so \texttt{NaN} is the
appropriate \texttt{\_FillValue} for Float64 throughout.

\section{Confidence $\times$ Importance: bucket summary}

The buckets below contain the proposed metadata. They are listed
once at the level of greatest generality (i.e.\ a rule that
applies to many variables is stated once and not repeated).

\subsection{\es{} -- \vc{}}\label{sec:es-vc}

These should be added by anyone re-publishing the dataset, in the
form below, before any consumer tool will accept the files.

\begin{enumerate}
  \item \textbf{Coordinate variables.} Add
    \texttt{lat(180)} and \texttt{lon(360)} as one-dimensional
    coordinate variables that share their names with the
    dimensions, with values
    \begin{quote}\ttfamily
      lat = -89.5, -88.5, \ldots, +89.5 \\
      lon = -179.5, -178.5, \ldots, +179.5
    \end{quote}
    and attributes
    \begin{quote}\ttfamily
      lat: units="degrees\_north", long\_name="Latitude",\\
      \phantom{lat:}\ standard\_name="latitude", axis="Y";\\
      lon: units="degrees\_east",  long\_name="Longitude",\\
      \phantom{lon:}\ standard\_name="longitude", axis="X".
    \end{quote}
    The mapping is established by the geographic test in
    \S\ref{sec:data}; the cell-centred placement is the standard
    convention for a $1^\circ \times 1^\circ$ binning.

  \item \textbf{Fill value on every Float64 variable.} The
    appropriate value is \texttt{\_FillValue = NaN}; observed in
    the data and consistent with the Float64 dtype.

  \item \textbf{Units on count variables.}
    \texttt{day\_pixel\_count}, \texttt{night\_pixel\_count},
    \texttt{day\_as\_pixel\_count}, \texttt{night\_as\_pixel\_count},
    \texttt{day\_at\_pixel\_count}, \texttt{night\_at\_pixel\_count}:
    \texttt{units="1"} (CF-style dimensionless count).

  \item \textbf{Units on SST sums.}
    \texttt{day\_sum\_SST}, \texttt{night\_sum\_SST}:
    \texttt{units="degree\_Celsius"} (inherited from the
    \texttt{SST\_Orbits/SST\_In} unit). The corresponding
    \texttt{*\_sum\_SST\_squared}: \texttt{units="degree\_Celsius2"}.

  \item \textbf{Units on geographic gradient sums.}
    \texttt{*\_sum\_eastward\_gradient},
    \texttt{*\_sum\_northward\_gradient},
    \texttt{*\_sum\_magnitude\_gradient}:
    \texttt{units="degree\_Celsius km-1"}. Squared variants:
    \texttt{units="degree\_Celsius2 km-2"}. Inherited from
    \texttt{SST\_Orbits/eastward\_gradient} (declared as
    \texttt{C/km}; the substitution
    \texttt{C} $\to$ \texttt{degree\_Celsius} was already proposed
    in the Prompt~\#2 audit).

  \item \textbf{Units on swath-relative gradient sums.}
    \texttt{*\_sum\_grad\_as\_per\_km},
    \texttt{*\_sum\_grad\_at\_per\_km},
    \texttt{*\_sum\_grad\_mag\_per\_km}: the suffix
    \texttt{\_per\_km} explicitly states the unit
    (\texttt{units="degree\_Celsius km-1"}). Squared variants:
    \texttt{units="degree\_Celsius2 km-2"}.

  \item \textbf{Long names.} For every variable, a long name can be
    formed mechanically from the components; see
    Table~\ref{tab:longnames}.

  \item \textbf{Global attribute set (minimum).}
    \texttt{Conventions="CF-1.8"};
    \texttt{title="Monthly day/night MODIS Aqua SST and SST-gradient
    statistics on a $1^\circ\times 1^\circ$ global grid"};
    \texttt{cdm\_data\_type="Grid"}.

  \item \textbf{Time coverage.} For each file, derive
    \begin{quote}\ttfamily
      time\_coverage\_start = "YYYY-MM-01T00:00:00Z" \\
      time\_coverage\_end   = "YYYY-MM-DDT23:59:59Z"
    \end{quote}
    where \texttt{MM/YYYY} come from the filename and
    \texttt{DD} is the last day of \texttt{MM}.

  \item \textbf{Geospatial extent.}
    \texttt{geospatial\_lat\_min=-90}, \texttt{lat\_max=+90},
    \texttt{lon\_min=-180}, \texttt{lon\_max=+180},
    \texttt{geospatial\_lat\_resolution="1 degree"},
    \texttt{geospatial\_lon\_resolution="1 degree"}.
\end{enumerate}

\subsection{\es{} -- \gs{}}\label{sec:es-gs}

These are essential but rest on inference; a reader should treat
the proposed values as best guesses that the producer should
confirm before publication.

\begin{enumerate}
  \item \textbf{Time variable (best guess).} Add a scalar time
    coordinate
    \begin{quote}\ttfamily
      time = (days since 1970-01-01) at the centre of \\
      \phantom{time = }the month encoded in the filename \\
      time: units="days since 1970-01-01 00:00:00 UTC", \\
      \phantom{time:}\ standard\_name="time", axis="T", \\
      \phantom{time:}\ calendar="gregorian", \\
      \phantom{time:}\ bounds="time\_bnds"
    \end{quote}
    with \texttt{time\_bnds = [first\_day, last\_day]} of the
    month. Guess because the file itself does not declare a time
    representation; mid-month UTC is the convention used by most
    monthly satellite products but cannot be confirmed without the
    producer.

  \item \textbf{Day / night definition (best guess).}
    \texttt{comment} on every \texttt{day\_*} variable:
    ``Aqua ascending node ($\sim$13:30 local solar time)''; on
    every \texttt{night\_*} variable:
    ``Aqua descending node ($\sim$01:30 local solar time)''. This
    is the standard MODIS-Aqua day/night split; guess because no
    explicit threshold is recorded in either branch.

  \item \textbf{`as'/`at' as along-scan/along-track (best
    guess).} \texttt{comment} on every
    \texttt{*\_grad\_as\_per\_km} variable:
    ``Along-scan (across-track) component of the SST gradient,
    in the swath-relative frame.'' Analogous for
    \texttt{*\_grad\_at\_per\_km} with ``along-track.'' Guess
    because the file does not spell out the abbreviation, but
    the naming and the std-dev parity with \texttt{eastward}/
    \texttt{northward} make this the only natural reading.

  \item \textbf{cell\_methods.}
    \begin{quote}\ttfamily
      *\_pixel\_count: cell\_methods="area: count time: sum" \\
      *\_sum\_SST:    cell\_methods="area: sum time: sum" \\
      *\_sum\_SST\_squared: cell\_methods="area: sum time: sum" \\
      *\_sum\_*\_gradient[\_squared]: cell\_methods="area: sum time: sum"
    \end{quote}
    This expresses the bookkeeping: each cell carries
    sums across the month and across all valid pixels in the
    cell, so that an aggregator can recompute means and
    variances. Guess in form (the CF \texttt{cell\_methods}
    syntax is not entirely unambiguous for accumulated sums)
    but the substance is firmly inferable from the variable
    naming.

  \item \textbf{Coordinates attribute.}
    \texttt{coordinates="time lat lon"} on every data variable,
    so that a CF-aware reader links the field to its axes
    (\texttt{time} is a scalar but its addition is required by
    CF when \texttt{time\_coverage\_*} is in the globals).

  \item \textbf{Per-variable computational notes for users.}
    Add as \texttt{comment} on each variable, the divisor that
    converts the sum to a mean (this is the most likely thing a
    user will get wrong):
    \begin{itemize}[topsep=0pt]
      \item \texttt{*\_sum\_SST} mean: divide by
        \texttt{*\_pixel\_count}.
      \item \texttt{*\_sum\_eastward\_gradient},
        \texttt{*\_sum\_northward\_gradient},
        \texttt{*\_sum\_magnitude\_gradient}: divide by
        \texttt{*\_pixel\_count}.
      \item \texttt{*\_sum\_grad\_as\_per\_km}: divide by
        \texttt{*\_as\_pixel\_count}.
      \item \texttt{*\_sum\_grad\_at\_per\_km}: divide by
        \texttt{*\_at\_pixel\_count}.
      \item \texttt{*\_sum\_grad\_mag\_per\_km}: divide by
        \texttt{*\_at\_pixel\_count} (verified empirically; see
        \S\ref{sec:data}).
    \end{itemize}
    The matching rule for the \texttt{\_squared} sums is the
    same divisor.
\end{enumerate}

\subsection{\es{} -- \nc{}}\label{sec:es-nc}

The audit would not feel comfortable proposing a value for these,
but they are essential for proper re-use:

\begin{itemize}
  \item \textbf{date\_created} of each file (per-file).
    The producer's filesystem and processing log have this; the
    file itself does not.
  \item \textbf{history.} The processing history (script,
    git commit / version, system, dependencies). Critical for
    reproducibility but absent.
  \item \textbf{Specific algorithm used to compute the
    upstream pixel-level gradients}, recorded as
    \texttt{comment} or \texttt{source}.
\end{itemize}

\subsection{\nh{} -- \vc{}}\label{sec:nh-vc}

\begin{enumerate}
  \item \textbf{ACDD provenance and contacts (inherited).} The
    sibling branch \texttt{SST\_Orbits/} declares these
    consistently across files. Reasonable to copy:
    \begin{quote}\ttfamily
      institution="University of Rhode Island - \\
      \phantom{institution=}Graduate School of Oceanography" \\
      creator\_name="Peter Cornillon" \\
      creator\_email="pcornillon@gso.uri.edu" \\
      creator\_url="http://www.sstfronts.org" \\
      project="SST Fronts" \\
      publisher\_name="SST Fronts" \\
      publisher\_institution=\textit{same as institution}
    \end{quote}
  \item \textbf{Source and lineage}.
    \begin{quote}\ttfamily
      source="MODIS Aqua L2 SST (URI reprocessing of NASA \\
      \phantom{source=}OceanColor L2 granules); the per-orbit \\
      \phantom{source=}files are at \dots/SST\_Orbits/." \\
      processing\_level="L4"   (monthly grid is L4) \\
      platform="Aqua" \\
      instrument="MODIS"
    \end{quote}

  \item \textbf{Conventions on the squared sums.} CF doesn't have
    a standard\_name for ``sum of squared SST'', so propose
    \begin{quote}\ttfamily
      *\_sum\_SST\_squared:\\
      \phantom{xx}long\_name="Sum of squared SST contributions" \\
      \phantom{xx}cell\_methods="area: sum\_of\_squares time: sum" \\
      \phantom{xx}comment="Use $\bar{x^2}-\bar{x}^2$ to recover variance."
    \end{quote}
    Same template for the gradient-squared sums.

  \item \textbf{Bounds variables for the coordinate axes.}
    \texttt{lat\_bnds(180,2)} = $[[-90,-89],\ldots,[89,90]]$;
    \texttt{lon\_bnds(360,2)} = $[[-180,-179],\ldots,[179,180]]$;
    referenced as \texttt{lat:bounds="lat\_bnds"} and
    \texttt{lon:bounds="lon\_bnds"}.
\end{enumerate}

\subsection{\nh{} -- \gs{}}\label{sec:nh-gs}

\begin{enumerate}
  \item \textbf{standard\_name on the gradient sums.} CF has
    \texttt{eastward\_temperature\_gradient} and
    \texttt{northward\_temperature\_gradient} for the components;
    the sums in this file are not themselves temperature
    gradients, so a strict CF reading would invent
    \texttt{standard\_name} extensions. Best to omit
    \texttt{standard\_name} entirely on the \texttt{sum\_*}
    variables and use \texttt{long\_name} +
    \texttt{cell\_methods} to carry the meaning.
  \item \textbf{ACDD keywords.} \texttt{keywords="EARTH SCIENCE >
    OCEANS > OCEAN TEMPERATURE > SEA SURFACE TEMPERATURE; EARTH
    SCIENCE > OCEANS > OCEAN TEMPERATURE > SST FRONTS; EARTH
    SCIENCE > OCEANS > OCEAN TEMPERATURE > SST GRADIENTS"} (GCMD
    style).
    \texttt{keywords\_vocabulary="GCMD Earth Science Keywords"}.
  \item \textbf{License and acknowledgement.} Copy verbatim from
    \texttt{SST\_Orbits/}: ``These data are openly available
    \dots'' Guess because the gradient files may be governed by
    a different licence than the orbit files; producer should
    confirm.
  \item \textbf{cf\_role / featureType.} Most likely
    \texttt{featureType="grid"}; not strictly needed since the
    file is implicitly a regular grid, but a CF reader may
    require it.
  \item \textbf{references.} \url{http://www.sstfronts.org/} as in
    \texttt{SST\_Orbits/}. Guess because the publication
    reference for the gradient-statistics derived product
    specifically may differ.
\end{enumerate}

\subsection{\nh{} -- \nc{}}\label{sec:nh-nc}

These are interesting attributes that one would normally want, but
that I do not have evidence to propose:

\begin{itemize}
  \item Quality flags that may have been applied at the
    aggregation step (e.g., minimum number of valid pixels per
    cell, exclusion of cloud-edge pixels). Are any cells
    deliberately suppressed?
  \item Exact relationship between this file and the related
    \texttt{gradients\_by\_period\_5day/} (which is 403 from the
    public OPeNDAP endpoint). Likely a 5-day analog of the
    monthly product, but I have no direct evidence.
  \item A \texttt{coverage\_content\_type} (ACDD): is this
    ``physicalMeasurement'', ``modelResult'',
    ``referenceInformation''? The sums are observational, the
    gridding is computational, the choice is not obvious.
  \item Whether the gradient computation was carried out on
    \texttt{SST\_In} or on the smoothed \texttt{regridded\_sst}
    in the orbit files. The orbit-file
    \texttt{eastward\_gradient} attributes don't say; this
    monthly file inherits whatever choice was made upstream.
\end{itemize}

\section{Per-variable proposed long\_name and units}\label{sec:longnames}

The names below follow the structure proposed in
\S\ref{sec:es-vc}. Generated mechanically from the variable name
grammar.

\footnotesize
\begin{longtable}{@{}L{5.9cm} L{3.4cm} L{6.9cm}@{}}
\caption{Proposed COARDS-essential metadata for every variable in a
\texttt{gradients\_by\_period} file. \emph{All} variables are
\texttt{Float64} on \texttt{(lon=360, lat=180)} with
\texttt{\_FillValue=NaN}. Confidence is \vc{} for everything in
this table; \texttt{long\_name} is generated mechanically from the
naming grammar.}\label{tab:longnames}\\
\toprule
\textbf{Variable} & \textbf{units} & \textbf{long\_name} \\
\midrule
\endfirsthead
\toprule
\textbf{Variable} & \textbf{units} & \textbf{long\_name} \\
\midrule
\endhead

\texttt{day\_pixel\_count}      & \texttt{1} & Number of valid daytime SST pixels in the cell over the month \\
\texttt{night\_pixel\_count}    & \texttt{1} & Number of valid night-time SST pixels in the cell over the month \\
\texttt{day\_as\_pixel\_count}  & \texttt{1} & Number of pixels with a valid along-scan gradient (daytime) \\
\texttt{night\_as\_pixel\_count}& \texttt{1} & Number of pixels with a valid along-scan gradient (night-time) \\
\texttt{day\_at\_pixel\_count}  & \texttt{1} & Number of pixels with a valid along-track gradient (daytime) \\
\texttt{night\_at\_pixel\_count}& \texttt{1} & Number of pixels with a valid along-track gradient (night-time) \\
\midrule
\texttt{day\_sum\_SST}                  & \texttt{degree\_Celsius}  & Sum of daytime SST contributions in the cell \\
\texttt{day\_sum\_SST\_squared}         & \texttt{degree\_Celsius2} & Sum of squared daytime SST contributions in the cell \\
\texttt{night\_sum\_SST}                & \texttt{degree\_Celsius}  & Sum of night-time SST contributions in the cell \\
\texttt{night\_sum\_SST\_squared}       & \texttt{degree\_Celsius2} & Sum of squared night-time SST contributions in the cell \\
\midrule
\texttt{day\_sum\_eastward\_gradient}            & \texttt{degree\_Celsius km-1} & Sum of daytime eastward SST gradient \\
\texttt{day\_sum\_eastward\_gradient\_squared}   & \texttt{degree\_Celsius2 km-2} & Sum of squared daytime eastward SST gradient \\
\texttt{day\_sum\_northward\_gradient}           & \texttt{degree\_Celsius km-1} & Sum of daytime northward SST gradient \\
\texttt{day\_sum\_northward\_gradient\_squared}  & \texttt{degree\_Celsius2 km-2} & Sum of squared daytime northward SST gradient \\
\texttt{day\_sum\_magnitude\_gradient}           & \texttt{degree\_Celsius km-1} & Sum of daytime SST gradient magnitude (geographic frame) \\
\texttt{day\_sum\_magnitude\_gradient\_squared}  & \texttt{degree\_Celsius2 km-2} & Sum of squared daytime SST gradient magnitude (geographic) \\
\texttt{night\_sum\_eastward\_gradient}          & \texttt{degree\_Celsius km-1} & Sum of night-time eastward SST gradient \\
\texttt{night\_sum\_eastward\_gradient\_squared} & \texttt{degree\_Celsius2 km-2} & Sum of squared night-time eastward SST gradient \\
\texttt{night\_sum\_northward\_gradient}         & \texttt{degree\_Celsius km-1} & Sum of night-time northward SST gradient \\
\texttt{night\_sum\_northward\_gradient\_squared}& \texttt{degree\_Celsius2 km-2} & Sum of squared night-time northward SST gradient \\
\texttt{night\_sum\_magnitude\_gradient}         & \texttt{degree\_Celsius km-1} & Sum of night-time SST gradient magnitude (geographic) \\
\texttt{night\_sum\_magnitude\_gradient\_squared}& \texttt{degree\_Celsius2 km-2} & Sum of squared night-time SST gradient magnitude (geographic) \\
\midrule
\texttt{day\_sum\_grad\_as\_per\_km}             & \texttt{degree\_Celsius km-1} & Sum of daytime along-scan SST gradient (swath-relative) \\
\texttt{day\_sum\_grad\_as\_per\_km\_squared}    & \texttt{degree\_Celsius2 km-2} & Sum of squared daytime along-scan SST gradient \\
\texttt{day\_sum\_grad\_at\_per\_km}             & \texttt{degree\_Celsius km-1} & Sum of daytime along-track SST gradient (swath-relative) \\
\texttt{day\_sum\_grad\_at\_per\_km\_squared}    & \texttt{degree\_Celsius2 km-2} & Sum of squared daytime along-track SST gradient \\
\texttt{day\_sum\_grad\_mag\_per\_km}            & \texttt{degree\_Celsius km-1} & Sum of daytime SST gradient magnitude (swath-relative frame) \\
\texttt{day\_sum\_grad\_mag\_per\_km\_squared}   & \texttt{degree\_Celsius2 km-2} & Sum of squared daytime SST gradient magnitude (swath) \\
\texttt{night\_sum\_grad\_as\_per\_km}           & \texttt{degree\_Celsius km-1} & Sum of night-time along-scan SST gradient \\
\texttt{night\_sum\_grad\_as\_per\_km\_squared}  & \texttt{degree\_Celsius2 km-2} & Sum of squared night-time along-scan SST gradient \\
\texttt{night\_sum\_grad\_at\_per\_km}           & \texttt{degree\_Celsius km-1} & Sum of night-time along-track SST gradient \\
\texttt{night\_sum\_grad\_at\_per\_km\_squared}  & \texttt{degree\_Celsius2 km-2} & Sum of squared night-time along-track SST gradient \\
\texttt{night\_sum\_grad\_mag\_per\_km}          & \texttt{degree\_Celsius km-1} & Sum of night-time SST gradient magnitude (swath) \\
\texttt{night\_sum\_grad\_mag\_per\_km\_squared} & \texttt{degree\_Celsius2 km-2} & Sum of squared night-time SST gradient magnitude (swath) \\
\bottomrule
\end{longtable}
\normalsize

\section{Cross-tabular summary}

\begin{center}
\renewcommand{\arraystretch}{1.15}
\begin{tabular}{@{}L{2.8cm} L{6.0cm} L{6.0cm}@{}}
\toprule
& \es{} & \nh{} \\
\midrule
\vc{}
& Coordinate variables (\texttt{lat}, \texttt{lon}), \texttt{\_FillValue=NaN}, all per-variable units and long\_names, time\_coverage and geospatial globals, Conventions/title/cdm\_data\_type.
& Inherited provenance from \texttt{SST\_Orbits/} (institution, creator\_*, project, publisher\_*); \texttt{source}, \texttt{platform}, \texttt{instrument}, \texttt{processing\_level}; coordinate \texttt{bounds} variables. \\
\midrule
\gs{}
& Scalar \texttt{time} variable (mid-month), day/night node \texttt{comment}s, as/at = along-scan/along-track \texttt{comment}s, \texttt{cell\_methods}, \texttt{coordinates}, per-variable divisor \texttt{comment}.
& \texttt{keywords} / \texttt{keywords\_vocabulary}, license, \texttt{featureType="grid"}, \texttt{references}. \\
\midrule
\nc{}
& \texttt{date\_created} per file, \texttt{history}, exact upstream-gradient algorithm.
& \texttt{coverage\_content\_type}, downstream relation to \texttt{gradients\_by\_period\_5day/}, choice of \texttt{SST\_In} vs.\ \texttt{regridded\_sst} as the upstream SST field, any minimum-pixel cell suppression. \\
\bottomrule
\end{tabular}
\end{center}

\section{Summary}

The \texttt{gradients\_by\_period/} branch is the harder of the
two readable branches: nothing is declared in the file, and a
recipient has to derive everything from the variable-name grammar,
from the data values, and from the sibling branch. The mechanical
inferences (\S\ref{sec:es-vc}) are not in dispute: the dimension
names \texttt{lat}/\texttt{lon} with sizes 180/360 plus the
land/ocean test pin the geographic grid; the suffix
\texttt{\_per\_km} and the sibling branch pin the units; the
Float64 dtype plus observed NaNs pin the fill value; the
\texttt{day\_}/\texttt{night\_} prefix plus the sibling branch
pin the temporal decomposition. The guesses (\S\ref{sec:es-gs})
are the points at which the writer's intentions become invisible
to the reader: which divisor to use for which sum, whether
\texttt{as}/\texttt{at} stand for along-scan/along-track, what
mid-month time stamp to record. The unguessable items
(\S\ref{sec:es-nc}, \S\ref{sec:nh-nc}) are essentially the
\emph{history} of the file --- when, by whom, by what code, with
which upstream choices --- and no amount of structural inference
recovers those.

If \texttt{grep-dap} ever gets to the point of \emph{writing}
inferred attributes back into a CF-compliant copy of files like
these, the rule of thumb that the audit suggests is: write
everything in \es{}--\vc{} silently, write everything in
\es{}--\gs{} but tag the attribute (e.g.\ via a
\texttt{comment} prefix \texttt{"[inferred]"}), and refuse to
write \es{}--\nc{} entries at all rather than fabricate them.

\section{Caveats}

\begin{itemize}
  \item Only the day-side of one month (\texttt{07\_2020}) and
    one month of contrasting season (\texttt{01\_2020}) were
    used for the sign / divisor / orientation tests. The
    same probe on a different month, or on a night-side
    variable, could reveal exceptions.
  \item All proposals are written against COARDS-recognizing
    consumers; CF-1.8 is closer to current practice and is the
    convention proposed for the \texttt{Conventions} attribute.
  \item The companion branch \texttt{gradients\_by\_period\_5day/}
    returns Hyrax 403; if its filename pattern matches
    (\texttt{augmented\_5day\_stats\_*}) the same inferences
    apply, but the schema cannot be verified without read access.
\end{itemize}

\end{document}
